% Separate post-CP15 candidate. A fibre bound, not an effective Hurwitz atlas.
\begin{proposition}[Proved: a quantitative fibre bound for fixed formal curve covers]
\label{D:fixed-formal-curve-fibre-bound}
Let $k$ be an algebraically closed field, let $E/k$ be a smooth proper
connected curve of genus $h$, and let $B\subset E(k)$ be finite. Fix
an integer $n\ge1$ and, at each $b\in B$, a finite free
$\widehat{\mathcal O}_{E,b}$-algebra $A_b$ of rank $n$ which is a
product of complete discrete valuation rings with residue field $k$
and whose generic algebra is etale. Write $d_b$ for its discriminant
exponent and put $D=\sum_{b\in B}d_b$.
Then the number of isomorphism classes over $E$ of finite generically
etale maps $f:C\to E$ of degree $n$, with $C/k$ smooth proper connected, which
are etale over $E\setminus B$ and satisfy
\[
       \prod_{x\in f^{-1}(b)}\widehat{\mathcal O}_{C,x}
          \simeq A_b
       \quad\text{as }\widehat{\mathcal O}_{E,b}\text{-algebras}
                                                               \tag{D.FC.1}
\]
for every $b$, is at most
\[
                        (n!)^{2h}\,n^{2D}.             \tag{D.FC.2}
\]
For $n=1$ the right side is understood as $1$. In particular, for
fixed $E$ and $D\le cn$, this is $n^{O_{h,c}(n)}$.
No Galois hypothesis, restriction on $p\mid n!$, lower bound on $n$
relative to the source genus, or finiteness hypothesis on $k$ is imposed.
The formal algebras are fixed in full, including their tame and
unramified factors.
\end{proposition}
\begin{proof}
The argument supplies a common characteristic-zero discriminant for
all the covers in one formal-data fibre. The closely related
\citet[Proposition 4.2, pp.\ 4--6]{DeJongKatzCurvesB} treats maps to
$\mathbb P^1$ with the same discriminant and formally isomorphic wild
ramification, under its hypotheses $g(C)\ge2$ and
$\deg(f)\ge2g(C)-1$. Its proof sketches the formal-patching step.
Here we fix every local factor, choose its lift directly, and give
that patching step explicitly. The resulting characteristic-zero
discriminant need not be reduced; this is why the permutation-defect
count below is useful.

\emph{The characteristic-zero count.}
First suppose $k$ has characteristic zero. For a cover in the
statement its discriminant divisor is the fixed divisor
$\Delta=\sum_b d_b[b]$. Let $b_1,\ldots,b_s$ be its support.
The fundamental group of $E\setminus\operatorname{Supp}(\Delta)$
has generators
\[
 x_1,y_1,\ldots,x_h,y_h,\sigma_1,\ldots,\sigma_s,
 \qquad \prod_i[x_i,y_i]\prod_j\sigma_j=1,             \tag{D.FC.3}
\]
where $\sigma_j$ generates inertia at $b_j$; this is the
characteristic-zero case, with its full proof, of
\citet[Expos\'e XIII, Corollary 2.12, pp.\ 392--394]{SGA1}.
That proof also handles arbitrary algebraically closed fields of
characteristic zero by change of algebraically closed field, so no
embedding of our entire field into $\mathbb C$ is assumed.

After labelling a geometric fibre, the cover gives a permutation
representation in $S_n$. If the inertia permutation at $b_j$ has
cycle lengths $e_1,\ldots,e_a$, its discriminant multiplicity is
\[
                    d_{b_j}=\sum_{i=1}^a(e_i-1)=n-a.   \tag{D.FC.4}
\]
Every permutation with index $n-a=d_{b_j}$ is a product of exactly
$d_{b_j}$ transpositions: decompose each of its cycles in this way.
There are at most $\binom n2^{d_{b_j}}\le n^{2d_{b_j}}$
such products. The $2h$ handle generators have at most $(n!)^{2h}$
possible images. Multiplying these estimates, and discarding both
the relation and transitivity when making an upper bound, gives
(D.FC.2). Quotienting labelled representations by conjugacy can
only decrease their number. The same bound therefore counts even
possibly disconnected smooth covers with this discriminant.

\emph{One lift of the target and of each prescribed local algebra.}
Assume now $\operatorname{char}(k)=p>0$, and put $R=W(k)$,
$K=\operatorname{Frac}(R)$. The ring $R$ is a complete discrete
valuation ring with residue field $k$. Choose a smooth proper lift
$\mathcal E/R$ of $E$. Existence follows from the unobstructed
deformation theory of smooth proper curves, followed by
algebraization; an ample line bundle lifts since $H^2(E,\mathcal O_E)=0$.
Equivalently one can use the smoothness statements in
\citet[Tags 0DZQ and 0DZT, Lemmas 93.16.4 and 109.15.1--2]{StacksFormalCurveB}.
This choice is made once, independently of $f$.
Lift the finitely many points $b$ to disjoint sections of
$\mathcal E/R$. Smoothness, henselian lifting and completeness
allow the choices, and their formal neighbourhoods have coordinates
$\widehat{\mathcal E}_{b_R}\simeq\operatorname{Spf}(R[[t_b]])$
lifting the chosen coordinates $\widehat{\mathcal O}_{E,b}=k[[t_b]]$.

For a nontrivial factor of $A_b$, choose a uniformizer and its
monic Eisenstein polynomial over $k[[t_b]]$. Lift every coefficient
to $R[[t_b]]$, still making all nonleading coefficients divisible
by $t_b$ and the constant coefficient divided by $t_b$ a unit.
If desired, Proposition~\ref{D:local-Eisenstein-linear-atlas}
first replaces the polynomial by a finite one representing the
same local extension. The lifted algebra is finite free over
$R[[t_b]]$. It is isomorphic to $R[[s]]$: writing the equation as
$s^e+t_b U(t_b,s)=0$, the constant $U(0,0)$ is a unit, so formal
implicit substitution solves uniquely for $t_b$ as a series in
$s$. The $t_b$- and $s$-adic topologies are cofinal. A degree-one
factor is lifted by the identity algebra. Thus we have fixed
algebras
\[
        \mathcal A_b=\prod_{j}R[[s_{bj}]]
           \quad\text{finite free of rank }n\text{ over }R[[t_b]],
        \qquad \mathcal A_b/p\mathcal A_b\simeq A_b.   \tag{D.FC.5}
\]
These choices need not preserve any local Galois group.

Let $\delta_b(t_b)$ be their trace discriminant. It reduces to
$t_b^{d_b}$ times a unit modulo $p$. In particular
$\mathcal A_b/p^m\mathcal A_b$ becomes finite etale after inverting
$t_b$, for every $m\ge1$: its discriminant is a unit in
$(R/p^m)((t_b))$. Weierstrass preparation also shows that
$R[[t_b]]/(\delta_b)$ is finite free over $R$ of rank $d_b$,
with the zero algebra allowed when $d_b=0$.

\emph{Patching at every nilpotent level.}
Fix any special-fibre cover $f:C\to E$ counted in the statement,
and fix isomorphisms (D.FC.1). Put $R_m=R/p^mR$,
$\mathcal E_m=\mathcal E\times_R R_m$, and let $U_m$ be the
complement of the chosen sections. The given finite etale cover
over $U_1$ has compatible finite etale lifts over all $U_m$.
Moreover every prescribed isomorphism on the reduction lifts
uniquely. This is the equivalence of finite-etale categories for
a nilpotent thickening, not an assertion about a characteristic-zero
punctured disc; see \citet[Tag 0BQB, Lemma 58.8.3]{StacksFormalCurveB}.

On each formal disc use the fixed finite free algebra
$\mathcal A_b/p^m\mathcal A_b$. On its puncture it and the
restriction of the lifted cover of $U_m$ are finite etale.
The chosen special-fibre identification therefore has a unique
lift to an isomorphism of these two algebras on
$R_m((t_b))$. The uniqueness makes these overlap isomorphisms
compatible in $m$.

For completeness, this datum really gives a finite locally free
algebra on $\mathcal E_m$, even when $R_m$ is nonreduced.
Near a branch section choose an affine neighbourhood on which
the section is cut out by one nonzerodivisor. Its complement is
already in $U_m$; its completion is the formal disc just used.
Beauville--Laszlo gluing for this nonzerodivisor gives an
equivalence of finite-projective modules between the affine
neighbourhood and its open/completed gluing data. The exact
statements, with complete proofs, are
\citet[Tag 0BP2 and Tag 0BNI, Theorem 15.92.16 and
Lemmas 15.92.18--19]{StacksFormalCurveB}.
Multiplication and unit descend by full faithfulness, because
tensor products have the corresponding base changes. Their
commutativity, associativity and unit identities can be checked
on the two pieces and consequently hold after gluing. The
resulting module has rank $n$ on both pieces and hence rank $n$
on the neighbourhood. Different branch neighbourhoods meet
inside $U_m$, so the resulting algebras glue there as well.
This constructs compatible finite locally free maps
\[
                     f_m:\mathcal C_m\longrightarrow\mathcal E_m
                                                               \tag{D.FC.6}
\]
with special fibre $f$ and the specified completed local algebras.

Grothendieck existence for finite covers of a proper scheme
algebraizes (D.FC.6) to a finite map
$\mathcal C\to\mathcal E$, proper over $R$; the finite-algebra
form is \citet[Tag 09ZT, Lemma 30.28.2]{StacksFormalCurveB}.
It is finite locally free of rank $n$. Indeed its coherent algebra
is locally free on the formal completion along the closed fibre,
so it is locally free at that fibre by faithful flatness of local
completion. Its nonfree locus is closed and proper over
$\operatorname{Spec}(R)$. Any nonempty such locus meets the closed
fibre, so it is empty. The source is therefore flat over $R$.
Its special fibre is the smooth curve $C$, which makes the
source smooth near that fibre. The nonsmooth locus is again
closed and proper over $R$, and the same argument proves that
$\mathcal C/R$ is smooth everywhere. We have thus constructed
a smooth proper lift of every $f$, using the same local lifts
(D.FC.5), with no bound on $n$ relative to $g(C)$.

\emph{One common global discriminant.}
The trace determinant of $\mathcal C\to\mathcal E$ defines an
effective Cartier divisor $\mathcal\Delta$ on $\mathcal E$.
It has no vertical component, since its reduction is the
nonzero discriminant of the generically etale map $f$.
It is proper and has finite fibres over $R$, hence is finite;
the absence of a vertical component makes it $R$-flat.
Its degree is therefore
\[
                     \deg_R(\mathcal\Delta)=\sum_b d_b=D.         \tag{D.FC.7}
\]
At every nilpotent level its ideal is the unit ideal on $U_m$
and is the fixed ideal $(\delta_b(t_b))$ on the formal disc at
each $b$. The full faithfulness in the preceding gluing step
therefore makes this ideal independent of $f$ at every level.
Full faithfulness of coherent formal algebraization then makes
$\mathcal\Delta$ itself independent of $f$.
If no special-fibre cover exists there is nothing to prove;
otherwise any one of the constructed covers supplies this common
algebraic divisor. In particular we do not assume in advance
that arbitrary formal divisors algebraize on the chosen target.

\emph{Geometric generic fibres distinguish special covers.}
Choose one constructed lift for each special-fibre isomorphism
class. Suppose two such lifts become isomorphic over
$\mathcal E_{\overline K}$, where $\overline K$ is an algebraic
closure of $K$. This isomorphism is defined over a finite
extension $K'/K$. Let $R'$ be the integral closure of $R$ in
$K'$; it is a finite complete discrete valuation ring, and its
residue field is still $k$. Both source models after base change
to $R'$ are smooth, hence regular and normal, and finite locally
free over $\mathcal E_{R'}$. Each is the normalization of the
target in its generic finite algebra: an element integral over
the target is integral over the finite source, and so belongs
to it by normality. Uniqueness of normalization extends the
generic isomorphism to the two entire finite covers over $R'$.
Reducing gives an isomorphism of the original covers over $E$.
Thus different special-fibre classes give different geometric
generic classes. This argument works with normal finite
algebras as well as with fields, so no hidden connectedness
claim after geometric generic base change is needed.

Finally, all geometric generic covers have target
$\mathcal E_{\overline K}$ of genus $h$ and the same discriminant
$\mathcal\Delta_{\overline K}$ of degree $D$. Apply the
characteristic-zero count (D.FC.3)--(D.FC.4). It proves
(D.FC.2), including the cases $B=\varnothing$ or $D=0$.
For fixed $h$ and $D\le cn$, use $n!\le n^n$ to obtain
the last assertion.
\end{proof}

\begin{remark}[Exact quantitative scope]
The proposition bounds the number of global covers in one fibre of
the map to complete local ramification data. In conjunction with
Proposition~\ref{D:local-Eisenstein-linear-atlas}, it separates a
local coefficient space of dimension $O(D)$ from a global fibre
of cardinality $n^{O(n+D)}$. It does not bound the degree of a
finite-type global gluing presentation, its equations, or the
complexity of an additional locus on which a cover extends over
a surface neighbourhood. A fibre-cardinality estimate alone
does not control the number of points in such a constrained
zero-dimensional locus. Those effective gluing and extension
conditions remain unproved here; in particular no new conclusion
about the general finite-generation questions D.1 or D.2 is made.
\end{remark}
